% 第 21 题答案辅助图：虚线仅标示原纸片剪线，不表示剪下后的拼合状态。
% 数字 1、2 标出所选的两块三角形底面；侧面配对与平移步骤见解析正文。
\begin{tikzpicture}[x=0.78cm,y=0.78cm,line cap=round,line join=round,
  outline/.style={draw=black,line width=0.65pt},
  cut/.style={draw=black,line width=0.55pt,densely dashed},
  font=\scriptsize]
  % 图 1：中点剪线将正三角形分成四个全等小三角形。
  \draw[outline] (1,1.732) -- (0,0) -- (2,0) -- cycle;
  \draw[cut]
    (0.5,0.866) -- (1.5,0.866)
    (0.5,0.866) -- (1,0)
    (1.5,0.866) -- (1,0);
  \node[below=2pt] at (1,-0.03) {图 1};

  % 图 2：中点剪线后，另两块角三角形再按内心—顶点连线切分。
  \begin{scope}[xshift=2.35cm]
    \draw[outline] (1,1.732) -- (0,0) -- (2,0) -- cycle;
    \draw[cut]
      (0.5,0.866) -- (1.5,0.866)
      (0.5,0.866) -- (1,0)
      (1.5,0.866) -- (1,0)
      (0.5,0.289) -- (0,0)
      (0.5,0.289) -- (0.5,0.866)
      (0.5,0.289) -- (1,0)
      (1.5,0.289) -- (2,0)
      (1.5,0.289) -- (1.5,0.866)
      (1.5,0.289) -- (1,0);
    \node at (1,0.58) {1};
    \node at (1,1.16) {2};
    \node[below=2pt] at (1,-0.03) {图 2};
  \end{scope}

  % 图 3：近似源图轮廓；两块底面与图 2 采用相同位置关系。
  \begin{scope}[xshift=4.75cm]
    \coordinate (A) at (0.75,1.75);
    \coordinate (B) at (0,0);
    \coordinate (C) at (4,0);
    \coordinate (Mab) at (0.375,0.875);
    \coordinate (Mac) at (2.375,0.875);
    \coordinate (Mbc) at (2,0);
    \coordinate (IL) at (0.553,0.365);
    \coordinate (IR) at (2.553,0.365);

    \draw[outline] (A) -- (B) -- (C) -- cycle;
    \draw[cut]
      (Mab) -- (Mac)
      (Mab) -- (Mbc)
      (Mac) -- (Mbc)
      (IL) -- (B)
      (IL) -- (Mab)
      (IL) -- (Mbc)
      (IR) -- (C)
      (IR) -- (Mac)
      (IR) -- (Mbc);
    \node at (1.583,0.583) {1};
    \node at (1.25,1.167) {2};
    \node[below=2pt] at (2,-0.03) {图 3};
  \end{scope}
\end{tikzpicture}
